\pagenumbering{arabic}

\chapter{INTRODUCTION}

\section{Background}

Electricity is an essential resource for households, businesses, and industries, and the increasing demand for electrical energy has created a need for more efficient and transparent methods of electricity management. Conventional electricity billing systems generally operate using postpaid models in which electricity is consumed first and payment is made later. Such systems provide limited awareness of real-time energy consumption and may involve delays in billing, manual intervention, and difficulties in monitoring individual consumption patterns.

The development of smart grid technologies has introduced digital methods for monitoring and managing electricity consumption. Smart grid systems can integrate communication, data processing, monitoring, and automated decision-making to provide improved visibility and control over electricity usage. Prepaid electricity systems are another approach that allows consumers to purchase electricity credit before consumption. By linking electricity usage with an available balance, prepaid systems can encourage consumers to monitor and manage their energy consumption more effectively.

Blockchain technology provides another opportunity for improving the transparency and reliability of prepaid electricity management. A blockchain-based system can maintain transaction records and execute predefined rules through smart contracts. In addition, machine learning techniques can be used to analyze electricity consumption patterns, detect abnormal behavior, and forecast future consumption.

The proposed \textbf{Decentralized Prepaid Smart Power Grid} combines these concepts into a software-based system for electricity management. The system uses simulated electricity telemetry, a FastAPI backend, a SQLite database, blockchain-based prepaid balance management, machine learning for anomaly detection and consumption forecasting, and a React-based monitoring dashboard. The system is intended to demonstrate how these technologies can work together to provide a more transparent, intelligent, and automated approach to prepaid electricity management.

The current project focuses on the software architecture and simulation of the proposed system. Physical electricity meters, ESP32 hardware, electrical relays, and actual power distribution infrastructure are outside the scope of the implementation. Instead, simulated telemetry is used to represent the data that would be generated by an edge device in a real-world deployment.


\section{Gap Identification}

Existing electricity management systems address different aspects of energy monitoring and billing, but these functionalities are often implemented separately. Smart metering systems mainly focus on collecting electricity consumption data, while prepaid systems focus on credit-based electricity access and billing. IoT-based energy monitoring systems provide remote access to consumption information, but they may depend on centralized application servers for data processing and account management.

Blockchain-based energy management systems can provide transparent and traceable transactions, but their integration with real-time consumption monitoring and automated electricity management is still an area for further exploration. Similarly, machine learning techniques can be used for detecting unusual energy consumption and predicting future demand, but these capabilities are not always integrated with prepaid electricity management and blockchain-based authorization.

This creates a gap for a unified system that combines prepaid electricity management, consumption monitoring, blockchain-based transaction handling, anomaly detection, and energy forecasting within a single software architecture. The proposed project addresses this gap by integrating these functionalities into a single decentralized prepaid smart power grid prototype.

The system provides simulated telemetry ingestion, prepaid balance management, anomaly detection, energy forecasting, blockchain-based authorization, and a web-based dashboard. By integrating these components, the project demonstrates a software architecture that can potentially be extended to physical smart-metering infrastructure in future work.


\section{Motivation}

The primary motivation for this project is the need for improved control, transparency, and awareness in electricity consumption and prepaid energy management. Conventional postpaid electricity systems do not always provide consumers with immediate information about their current consumption or remaining financial commitment. A prepaid mechanism can provide a more direct relationship between electricity consumption and available credit.

Another motivation is the increasing use of digital technologies for intelligent energy management. Blockchain can provide a transparent mechanism for managing prepaid balances and authorization, while machine learning can help identify unusual consumption patterns and estimate future energy usage. Combining these technologies in a single system provides an opportunity to explore a more intelligent and decentralized approach to electricity management.

The project is also motivated by the need for a practical software prototype that can demonstrate these concepts without requiring physical electrical infrastructure. By using simulated telemetry, the proposed system can be developed and tested in a controlled environment while maintaining an architecture that can potentially support hardware integration in future implementations.


\section{Objectives}

The main objective of this project is to design and develop a software-based \textbf{Decentralized Prepaid Smart Power Grid} that integrates prepaid electricity management, blockchain technology, machine learning, and real-time monitoring.

The specific objectives of the project are:

\begin{itemize}

    \item To design a software architecture for decentralized prepaid electricity management.

    \item To implement a prepaid electricity mechanism using blockchain and smart contracts.

    \item To provide a mechanism for recording and managing simulated electricity consumption telemetry.

    \item To implement automatic prepaid balance deduction based on simulated energy consumption.

    \item To implement user authorization and cutoff mechanisms using blockchain smart-contract functions.

    \item To develop an anomaly detection mechanism using the Isolation Forest machine-learning algorithm.

    \item To develop an electricity consumption forecasting mechanism using linear regression.

    \item To provide a web-based dashboard for monitoring electricity consumption, prepaid balance, forecasting results, and anomaly information.

    \item To integrate the backend, database, blockchain, machine-learning components, and frontend into a unified software system.

    \item To evaluate the proposed system using simulated electricity telemetry without requiring physical electrical hardware.

\end{itemize}